\documentclass[12pt, a4paper]{article}
\usepackage[T2A]{fontenc}
\usepackage[utf8]{inputenc}
\usepackage[russian]{babel}
\usepackage{amsmath, amssymb}
\usepackage{enumitem}
\usepackage{geometry}
\geometry{top=2cm, bottom=2cm, left=2.5cm, right=2cm}

\begin{document}

\section*{Георгий Игоревич Вольфсон <<Делимость с человеческим лицом>>}

\subsection*{Беседа пятая. Признаки делимости.}

\subsubsection*{Упражнения}

\begin{enumerate}[label=\arabic*., wide=0pt, leftmargin=*]

\item \textbf{Для числа 1147 найдите ближайшее к нему натуральное число, кратное 3.}

Ответ: 1146.

\item \textbf{Укажите все цифры, которыми можно заменить звёздочку так, чтобы число делилось на 9: \(\ast 45\).}

Ответ: 9.

\item \textbf{Замените звёздочки двумя одинаковыми цифрами так, чтобы число \(11\ast\ast\) делилось на 15.}

Ответ: 1155.

\item \textbf{Докажите, что \(3^{100}+1\) делится на 2.}

\textbf{Решение:}
Число 3 — нечётное. Произведение любых нечётных чисел нечётно, поэтому \(3^{100}\) — нечётное. Тогда \(3^{100}+1\) — сумма нечётного и 1 (нечётное+нечётное=чётное). Чётное число всегда делится на 2. 
Таким образом, \(3^{100}+1\) кратно 2.

\item \textbf{Вычеркните в числе 21 945 две цифры так, чтобы число делилось на 9.}

Ответ: 945.

\item \textbf{Докажите, что ребус \(\text{АБ} \times \text{БГ} = \text{ДДЕЕ}\) не имеет решений.}

\textbf{Решение:}
Правая часть делится на 11 (11 * 100 * Д + 11 * Е), левая часть состоит из двух множителей, каждый из которых на 11 не делится, так как 11 число простое, значит левая часть не 11 не делится, значит решений нет.

\item \textbf{Укажите все цифры, которыми можно заменить звёздочку так, чтобы число делилось на 8: \(31\,415\,9\ast 6\).}

Ответ: 3 и 7.

\item \textbf{К числу 10 слева и справа припишите по одной цифре, чтобы полученное число делилось на 72. Найдите все такие варианты.}

\textbf{Решение:}
Число имеет вид \(\overline{x10y}\). Делимость на 72 означает делимость на 8 и на 9.
\begin{itemize}
\item На 8: последние три цифры \(\overline{10y}\) должны делиться на 8. Перебираем \(y=0,\dots,9\): 100,101,...,109. Из них на 8 делится только 104 (\(y=4\)).
\item На 9: сумма цифр числа \(x+1+0+4 = x+5\) должна делиться на 9. \(x+5\) кратно 9 при \(x=4\) (4+5=9). \(x\) — цифра, подходит.
\end{itemize}
Единственное число: 4104. 
Ответ: 4104.

\item \textbf{Может ли число, составленное только из троек, делиться на число, составленное только из восьмёрок?}

\textbf{Решение:}
Число из одних троек: \(33\ldots3\) — нечётное (так как оканчивается на 3). Число из одних восьмёрок: \(88\ldots8\) — чётное (оканчивается на 8). Чётное число не может быть делителем нечётного, так как при делении нечётного на чётное не получается целого числа. Следовательно, не может. 
Ответ: нет.

\item \textbf{Может ли число, составленное только из восьмёрок, делиться на число, составленное только из троек?}

Ответ: да (например, 888 : 3 = 296).

\item \textbf{Можно ли составить из цифр 0, 1, 2, ..., 7 восьмизначное число, которое будет делиться на 9? Каждую цифру можно использовать ровно один раз.}

\textbf{Решение:}
Сумма всех цифр от 0 до 7: \(0+1+2+3+4+5+6+7 = 28\). Признак делимости на 9: число делится на 9, если сумма его цифр делится на 9. 28 не делится на 9 (28:9 = 3·9=27, остаток 1). Поскольку перестановка цифр не меняет их сумму, любое восьмизначное число, составленное из этих цифр, будет иметь сумму цифр 28, а значит, не будет кратно 9. 
Ответ: нет, нельзя.

\end{enumerate}

\end{document}